\section{The Latin Fathers before Dionysius}\label{sec:fathers-latin}

The Latin Church spoke of the angels long before it wrote a treatise on
them. Its first witnesses were apologists who met the gods of Rome as
demons and drove them out in the name of Christ; then bishops and
exegetes who found the angels in the Psalms, the Prophets, and the
Gospels; then a monk who carried to Gaul what the fathers of the Egyptian
desert had learned of the spirits in combat. Across three centuries they
teach with one voice that God made the angels good and that the devil and
his companions fell by their own will; that the demons deceive by
subtlety and speed, and yield to the name of Christ; that the holy angels
praise God without ceasing, carry the prayers of the faithful, and keep
each soul from its birth; and that the heavenly powers stand in ranks
whose names Scripture gives. On the manner of the first sin, the bodies
of spirits, the princes of the nations, and the time of the angels'
creation they speak in their own voices and divide. Augustine gathers
this witness, and the Latin schools received his teaching
(\S\ref{sec:fl-augustine}).

\sectionguard
\subsection{Tertullian}\label{sec:fl-tertullian}

Tertullian of Carthage (d.~about 220) speaks of the angels and the demons
as an apologist before the magistrates of Rome,
as a controversialist against Marcion and his disciples, and as a teacher
of prayer and patience, and in each office the spirits enter his argument.

\paragraph{The demons in the \emph{Apology}.} Tertullian begins from what
the pagans themselves concede. The philosophers acknowledge demons,
Socrates waited on one, the poets know them, the common people curse by
Satan ``as though from some instinctive soul-knowledge of him,'' and the
magicians are ``witnesses to the existence of both kinds of spirits.''
Scripture adds their origin: ``We are instructed, moreover, by our sacred
books how from certain angels, who fell of their own free-will, there
sprang a more wicked demon-brood, condemned of God along with the authors
of their race, and that chief we have referred to'' (\emph{Apol.}\ 22.3).
Their business is single: ``Their great business is the ruin of mankind.''
They lay diseases on the body and drive the soul into sudden excesses, and
their nature fits them for both: ``Their marvellous subtleness and
tenuity give them access to both parts of our nature''
(\latin{subtilitas et tenuitas}, 22.5). Their action is known only by its
effects, like the unseen blight that withers the grain in the flower.

Their speed explains how they counterfeit divinity:
\begin{quote}
Every spirit is possessed of wings. This is a common property of both
angels and demons. So they are everywhere in a single moment; the whole
world is as one place to them; all that is done over the whole extent of
it, it is as easy for them to know as to report. Their swiftness of motion
is taken for divinity, because their nature is unknown.
(\emph{Apol.}\ 22.8)
\end{quote}
\noindent The Latin is terse: \latin{Omnis spiritus ales est: hoc angeli
et daemones.} What the demons foretell they have borrowed. ``The purposes
of God, too, they took up of old from the lips of the prophets, even as
they spoke them; and they gather them still from their works, when they
hear them read aloud''; their weather-lore comes from ``their commerce
with the clouds.'' Their cures are a fraud: ``first of all, they make you
ill; then, to get a miracle out of it, they command the application of
remedies~\dots\ and straightway withdrawing hurtful influence, they are
supposed to have wrought a cure.''

In the next chapter Tertullian lays down a challenge to the courts. ``Let
a person be brought before your tribunals, who is plainly under
demoniacal possession. The wicked spirit, bidden to speak by a follower of
Christ, will as readily make the truthful confession that he is a demon,
as elsewhere he has falsely asserted that he is a god'' (23.4). The
Christian's power over them is the name of the Judge:
\begin{quote}
Why, all the authority and power we have over them is from our naming the
name of Christ, and recalling to their memory the woes with which God
threatens them at the hands of Christ as Judge, and which they expect one
day to overtake them. Fearing Christ in God, and God in Christ, they
become subject to the servants of God and Christ. So at our touch and
breathing, overwhelmed by the thought and realization of those judgment
fires, they leave at our command the bodies they have entered, unwilling,
and distressed, and before your very eyes put to an open shame.
(\emph{Apol.}\ 23.15--16)
\end{quote}
\noindent Their confession becomes a testimony for the Gospel: ``your very
gods kindle up faith in our Scriptures, they build up the confidence of
our hope.''

\paragraph{The angels who sinned with the daughters of men.} The
``authors of their race'' in \emph{Apol.}\ 22.3 are, for Tertullian, the
angels of Gen 6:2, ``The sons of God seeing the daughters of men, that
they were fair, took to themselves wives of all which they chose.'' He
names them ``those angels, to wit, who rushed from heaven on the daughters
of men,'' and teaches that they disclosed to men the working of metals,
the powers of herbs and enchantments, and ``every curious art, even to the
interpretation of the stars,'' with the ornaments and paints of women
(\emph{De cultu feminarum} I.2). The astrologers are condemned with ``those
angels, the deserters from God, the lovers of women,'' who discovered
their art: ``The city and Italy are interdicted to the astrologers, just
as heaven to their angels'' (\emph{De idololatria} 9). He reads the
Apostle's ``because of the angels'' (1~Cor 11:10) by the same history:
women are to be veiled ``because on account of `the daughters of men'
angels revolted from God'' (\emph{De oratione} 22). Cyprian and Lactantius
share this reading; Cassian answers it with the line of Seth
(\S\ref{sec:fl-cassian}), and Aquinas, following Augustine, reads the sons of
God as the sons of Seth (\ST{I q.~51 a.~3 ad~6}).

\paragraph{The devil made good (\emph{Adversus Marcionem} II.10).} Marcion
blamed the Creator for the devil who tempted Adam. Tertullian answers
that what God made was an angel, and that the devil made himself. The
malice and the slander of God in Paradise came ``most certainly not from
God, who made the angel good after the fashion of His good works. Indeed,
before he became the devil, he stands forth the wisest of creatures; and
wisdom is no evil.'' The proof is Ezekiel's lament over the prince of
Tyre, which he reads of the devil: ``this angel was both by creation good
and by choice corrupt.'' Glossing ``full of wisdom, perfect in beauty,''
he adds ``(this belongs to him as the highest of the angels, the
archangel, the wisest of all)''; and he shows that the lament cannot fit a
man, for none was born in God's paradise ``nor placed with a cherub upon
God's holy mountain, that is to say, in the heights of heaven, from which
the Lord testifies that Satan fell'' (Ezek 28:12--16; Luke 10:18).

The fall was an act of freedom God had given for glory: ``he was himself
as a spirit no less (than man) created, with the faculty of free-will. For
God would in nothing fail to endow a being who was to be next to Himself
with a liberty of this kind.'' He departed from his created condition
``through his own lusting after the wickedness which was spontaneously
conceived within him.'' Why was he not destroyed at once? Because God
``afforded room for a conflict, wherein man might crush his enemy with the
same freedom of his will as had made him succumb to him~\dots\ wherein
also the devil might receive a more bitter punishment, through being
vanquished by him whom he had previously injured.'' The Fourth Lateran
Council confesses the same in its creed: the devil and the other demons
were created good by nature and became evil of themselves (DS~800;
\S\ref{sec:faith}). Aquinas takes the covering cherub of Ezek 28:14 for the
highest angel who sinned, and holds it more probable that he was the
highest of all (\ST{I q.~63 a.~7} and ad~1; see \S\ref{sec:q63}).

\paragraph{Impatience and envy (\emph{De patientia} 5).} Treating patience,
Tertullian looks for the birth of its contrary in the adversary, and finds
it at the creation of man:
\begin{quote}
Therefore I detect the nativity of impatience in the devil himself, at
that very time when he impatiently bore that the Lord God subjected the
universal works which He had made to His own image, that is, to man. For
if he had endured (that), he would not have grieved; nor would he have
envied man if he had not grieved. Accordingly he deceived him, because he
had envied him; but he had envied because he had grieved: he had grieved
because, of course, he had not patiently borne. (\emph{De pat.}\ 5)
\end{quote}
\noindent The Latin has \latin{natales inpatientiae in ipso diabolo
deprehendo} (5.5). Whether malice or impatience came first in ``that
angel of perdition'' he declines to settle: the two ``grew up indivisible
in one paternal bosom.'' Impatience, conceived of the devil's seed, bore
anger and then Cain's murder. Aquinas determines that the first sin of the
angel was pride and that envy followed from it (\ST{I q.~63 a.~2}).

\paragraph{The bodies the angels assume (\emph{De carne Christi} 3 and 6).}
Against Marcion, who gave Christ a phantom flesh, Tertullian appeals to the
angels of the Old Testament, who ``have been changed into human form, and
have even borne about so veritable a body, that Abraham even washed their
feet, and Lot was rescued from the Sodomites by their hands; an angel,
moreover, wrestled with a man so strenuously with his body, that the
latter desired to be let loose'' (Gen 18--19; 32). If angels remained
angels while bearing true bodies, God can become man without ceasing to be
God. Their bodies came from nothing and returned to nothing, yet ``there
was solidity in their bodily substance'' (ch.~3).

Against Apelles, who gave Christ an unborn flesh on the pattern of the
angels, Tertullian draws the contrast between the missions. ``Never did
any angel descend for the purpose of being crucified, of tasting death,
and of rising again from the dead~\dots\ They had not come to die,
therefore they also (came not) to be born.'' He then states his own view
of the angelic nature and power:
\begin{quote}
It is plain that the angels bore a flesh which was not naturally their
own; their nature being of a spiritual substance, although in some sense
peculiar to themselves, corporeal; and yet they could be transfigured
into human shape, and for the time be able to appear and have intercourse
with men. Since, therefore, it has not been told us whence they obtained
their flesh, it remains for us not to doubt in our minds that a property
of angelic power is this, to assume to themselves bodily shape out of no
material substance. (\emph{De carne Christi} 6)
\end{quote}
\noindent The Latin makes the power a property of the angel:
\latin{proprium angelicae potestatis, ex nulla materia corpus sibi
sumere} (6.10). The angelic nature is \latin{substantiae spiritalis---etsi
corporis alicuius, sui tamen generis} (6.9), spiritual, yet a body of its
own kind: Tertullian, who holds even the soul to be corporeal (\emph{De
anima} 5--7), gives the angels a body of their own. Aquinas determines
that the angels have no bodies naturally united to them, and that the
bodies they assume are formed of air, ``condensing it by the Divine power
in so far as is needful'' (\ST{I q.~51 aa.~1--2}; a.~2 ad~3; see
\S\ref{sec:q51}).

\paragraph{The angels at prayer (\emph{De oratione}).} Expounding ``Hallowed
be thy name,'' Tertullian asks when the name of God is not holy, since it
sanctifies all things, ``He to whom that surrounding circle of angels
cease not to say, `Holy, holy, holy?' In like wise, therefore, we too,
candidates for angelhood, if we succeed in deserving it, begin even here
on earth to learn by heart that strain hereafter to be raised unto God,
and the function of future glory'' (\emph{De orat.}\ 3.3; Isa 6:3). The
Christians who pray are \latin{angelorum~\dots\ candidati}, learning now
the heavenly office they will keep for ever. When he rebukes those who sit
down as soon as prayer is ended, it is irreverence ``under the eye of the
living God, while the angel of prayer is still standing by'' (ch.~16). His
treatise closes on prayer as the praise of the whole creation: ``The
angels, likewise, all pray; every creature prays'' (ch.~29; \latin{Orant
etiam angeli omnes}). The Church still joins the angels in the thrice-holy
hymn (CCC~335).

\paragraph{Angels and demons about the soul (\emph{De anima}).} In the
treatise on the soul the good angels appear at birth and death. The pagans
assigned goddesses to every stage of gestation; ``We, on our part, believe
the angels to officiate herein for God'' (ch.~37). At death the soul, set
free, ``exults or it fears, according as it finds what lodging is prepared
for it, as soon as it sees the very angel's face, that arraigner of
souls'' (ch.~53). The demons attend the same thresholds. Among the heathen
the evil spirit claims the child at birth through the rites that
consecrate it to idols, and ``to all persons their genii are assigned,
which is only another name for demons'' (ch.~39). Under the cover of the
dead the demons counterfeit ghosts, and in exorcism a demon will claim to
be a kinsman of the possessed, a gladiator, or a god; ``and yet for all
that, the demon, after trying to circumvent the bystanders, is vanquished
by the pressure of divine grace, and sorely against his will confesses all
the truth'' (ch.~57).

\sectionguard
\subsection{Minucius Felix and Cyprian}\label{sec:fl-minucius-cyprian}

\paragraph{The \emph{Octavius}.} In the dialogue of Minucius Felix the
Christian Octavius answers the pagan appeal to auguries and oracles by
going to ``the very source of error and perverseness.'' He describes the
demons as fallen spirits:
\begin{quote}
There are some insincere and vagrant spirits degraded from their heavenly
vigour by earthly stains and lusts. Now these spirits, after having lost
the simplicity of their nature by being weighed down and immersed in
vices, for a solace of their calamity, cease not, now that they are ruined
themselves, to ruin others; and being depraved themselves, to infuse into
others the error of their depravity and being themselves alienated from
God, to separate others from God by the introduction of degraded
superstitions. (\emph{Octavius} 26)
\end{quote}
\noindent Chapter 27 sets out their works: consecrated under statues and images,
they ``lurk there, and by their afflatus attain the authority as of a
present deity,'' breathe into the prophets, ``animate the fibres of the
entrails, control the flights of birds, direct the lots.'' ``For they are
both deceived, and they deceive; inasmuch as they are both ignorant of the
simple truth, and for their own ruin they confess not that which they
know.'' Like Tertullian, Minucius rests the proof on exorcism. The demons
confess what they are ``as often as they are driven by us from bodies by
the torments of our words and by the fires of our prayers,'' and ``when
abjured by the only and true God, unwillingly the wretched beings shudder
in their bodies, and either at once leap forth, or vanish by degrees, as
the faith of the sufferer assists or the grace of the healer inspires''
(ch.~27).

\paragraph{Cyprian.} Cyprian of Carthage (d.~258) takes up the same
description almost word for word. The spirits worshipped in the idols
``are impure and wandering spirits, who, after having been steeped in
earthly vices, have departed from their celestial vigour by the contagion
of earth, and do not cease, when ruined themselves, to seek the ruin of
others'' (\emph{Quod idola dii non sint} 6). They lurk under statues and
``excite diseases to force them to worship of themselves''; adjured
through the true God, they ``at once yield and confess, and are
constrained to go out from the bodies possessed'' (ch.~7). To the
proconsul Demetrianus he makes the challenge Tertullian had made: ``Oh,
would you but hear and see them when they are adjured by us, and tortured
with spiritual scourges, and are ejected from the possessed bodies with
tortures of words, when howling and groaning at the voice of man and the
power of God, feeling the stripes and blows, they confess the judgment to
come!'' (\emph{Ad Demetrianum} 15).

Cyprian's own contribution is the account of the devil's fall by envy. In
the treatise \emph{On Jealousy and Envy} he traces that vice to its source
``at the very beginnings of the world'':
\begin{quote}
He who was sustained in angelic majesty, he who was accepted and beloved
of God, when he beheld man made in the image of God, broke forth into
jealousy with malevolent envy---not hurling down another by the instinct
of his jealousy before he himself was first hurled down by jealousy,
captive before he takes captive, ruined before he ruins others.
(\emph{De zelo et livore} 4)
\end{quote}
\noindent ``How great an evil is that, beloved brethren, whereby an angel
fell, whereby that lofty and illustrious grandeur could be defrauded and
overthrown, whereby he who deceived was himself deceived!'' From that
envy, as Wisdom says, death entered the world (Wis 2:24), and ``they who
are on his side imitate him.'' Cain's hatred of Abel, Esau's of Jacob, the
brothers' of Joseph, Saul's of David all descend from the first angel's
envy. With Tertullian, Cyprian also names the ``sinning and apostate
angels'' as the teachers of cosmetics, who ``forsook their heavenly
vigour'' (\emph{De habitu virginum} 14).

\sectionguard
\subsection{Lactantius}\label{sec:fl-lactantius}

Lactantius, the African teacher of rhetoric whom Constantine charged with
the education of his son Crispus, gave in the second book of the
\emph{Divine Institutes}, ``Of the Origin of Error,'' a full narrative of
his own on the origin of the devil and the demons.\footnote{Book and chapter numbers follow Brandt's edition (CSEL~19,
1890). The ANF translation numbers the chapters of Book~II cited here one
higher (II.9, 15--17) and the \emph{Epitome} chapters 22--23 as 27--28.}

\paragraph{The origin of the devil.} Before the world, Lactantius teaches,
God ``produced a Spirit like to Himself, who might be endowed with the
perfections of God the Father.'' He continues:
\begin{quote}
Then He made another being, in whom the disposition of the divine origin
did not remain. Therefore he was infected with his own envy as with
poison, and passed from good to evil; and at his own will, which had been
given to him by God unfettered, he acquired for himself a contrary name.
From which it appears that the source of all evils is envy. For he envied
his predecessor, who through his stedfastness is acceptable and dear to
God the Father. (\emph{Div.\ inst.}\ II.8.4--5)
\end{quote}
\noindent The Latin gives the principle in five words: \latin{cunctorum
malorum fontem esse liuorem}. This being ``who from good became evil by
his own act, is called by the Greeks diabolus: we call him accuser,
because he reports to God the faults to which he himself entices us''
(II.8.6). Lactantius here makes the devil's envy the envy of the first-begotten Son;
Tertullian, Cyprian, and Cassian place it on man, and the Church confesses
with all of them that the devil became evil by his own will (DS~800).

\paragraph{The angels sent to guard men.} Lactantius's account of the other
demons begins with a mission of the holy angels. When men had multiplied,
God, foreseeing that the devil, ``to whom from the beginning He had given
power over the earth,'' would corrupt or destroy them, ``sent angels for
the protection and improvement of the human race'' (\latin{misit angelos
ad tutelam cultumque generis humani}, II.14.1). Because He had given them
free will, He commanded them above all ``not to defile themselves with
contamination from the earth, and thus lose the dignity of their heavenly
nature.'' The ruler of the earth drew them, while they dwelt among men,
into intercourse with women. ``Then, not being admitted into heaven on
account of the sins into which they had plunged themselves, they fell to
the earth. Thus from angels the devil makes them to become his satellites
and attendants.'' Their offspring, ``neither angels nor men, but bearing a
kind of mixed nature,'' were not admitted into hell, as their fathers were
not into heaven: ``Thus there came to be two kinds of demons; one of
heaven, the other of the earth. The latter are the wicked spirits, the
authors of all the evils which are done, and the same devil is their
prince'' (II.14.3--5).

Their knowledge is real but bounded: ``They are acquainted, indeed, with
many future events, but not all, since it is not permitted them entirely
to know the counsel of God; and therefore they are accustomed to
accommodate their answers to ambiguous results'' (II.14.6). The poet who calls the
demons ``the guardians of mortal men'' has half the truth, ``because God
had sent them as guardians to the human race; but they themselves also,
though they are the destroyers of men, yet wish themselves to appear as
their guardians, that they themselves may be worshipped, and God may not
be worshipped'' (II.14.7--8). They cling to households under the name of
\latin{genii} and, being spirits, ``insinuate themselves into the bodies
of men'' (II.14).

\paragraph{The demons subject to the just.} They injure, Lactantius
continues, ``but those only by whom they are feared, whom the powerful and
lofty hand of God does not protect.'' They fear the worshippers of God,
``adjured by whose name they depart from the bodies of the possessed: for,
being lashed by their words as though by scourges, they not only confess
themselves to be demons, but even utter their own names~\dots\ because
they cannot speak falsely to God, by whom they are adjured, nor to the
righteous, by whose voice they are tortured'' (II.15).

\paragraph{The obedience of the holy angels.} Against the pagan pantheon he
sets the holy angels in their true place. The demons ``feigned that there
are many heavenly beings, and one king of all, Jupiter; because there are
many spirits of angels in heaven, and one Parent and Lord of all, God.''
The angels refuse divine honour:
\begin{quote}
nor do the angels, inasmuch as they are immortal, either suffer or wish
themselves to be called gods: for their one and only duty is to submit to
the will of God, and not to do anything at all except at His command. For
we say that the world is so governed by God, as a province is by its
ruler; and no one would say that his attendants are his sharers in the
administration of the province, although business is carried on by their
service.~\dots\ nor is there anything in the angels except the necessity
of obedience. Therefore they wish no honour to be paid to them, since all
their honour is in God. (\emph{Div.\ inst.}\ II.16)
\end{quote}
\noindent The image of the province and its officers gives the angels
their part in the government of the world; Aquinas gives its reason in the order of universal and particular powers
(\ST{I q.~110 a.~1}; \S\ref{sec:government}).

\paragraph{The \emph{Epitome}.} Lactantius's own abridgement gives the
story its sharpest form: the angels sent ``to instruct the race of men,
and to protect them from all evil'' were seduced, and ``being condemned by
the sentence of God, and cast forth on account of their sins, they lost
both the name and substance of angels. Thus, having become ministers of
the devil, that they might have a solace of their ruin, they betook
themselves to the ruining of men, for whose protection they had come''
(\emph{Epit.}\ 22.9--11). Cassian and the later Latin tradition place the fall
of the rebel angels with their leader, before the fall of man
(\S\ref{sec:fl-cassian}; \S\ref{sec:fall}); the mission of the holy angels as
guardians of the human race, which Lactantius sets at the head of his
narrative, is the teaching of the whole tradition (\ST{I q.~113 a.~1}).

\sectionguard
\subsection{Hilary of Poitiers}\label{sec:fl-hilary}

Hilary of Poitiers (d.~367), the defender of the Nicene faith in Gaul,
expounded the Psalms in a series of tractates in which the angels are a
constant presence.\footnote{Hilary's Latin follows Zingerle's edition,
CSEL~22 (1891), with its orthography and section numbers.}

\paragraph{A world full of angels.} On ``Blessed are they that search his
testimonies'' (Ps 118[119]:2), Hilary asks what God's testimonies are,
and answers that the whole of creation is under witnesses. The Apostle
charges Timothy ``before God and Christ Jesus and the elect angels'' (1~Tim
5:21), and so no one lives in an empty world: \latin{plena sunt enim omnia
diuinis testimoniis; et omne hoc, uacuum quod putatur, repletum angelis
dei nihilque est, quod non haec diuinorum ministeriorum frequentia
incolat}. All things are full of God's witnesses; all this space that is
thought empty is filled with the angels of God, and there is nothing that
this throng of divine ministers does not inhabit (\emph{Tract.\ in Ps.}\
118, Aleph 7). Men will not steal before a witness; how much more should a
Christian fear to sin, knowing himself surrounded on every side by the
witness of spiritual powers? \latin{metuimus adsistentes undique nobis choros
angelorum et plenum ministeriis caelestibus mundum?}---do we fear the
choirs of angels standing round us and a world full of heavenly
ministries? If the angels of the little ones see the Father daily (Matt
18:10), we may fear the witness of those whom we know both to remain
with us and daily to stand before God. The devil, too, is a witness,
\latin{puncto temporis omnem amplitudinem mundi huius obeuntem}, going
round the whole breadth of the world in a moment, to whom it is sweet
that we sin, that he may glory in the witness of our sins; for this is
his proper cunning, \latin{instigare ad peccandum et arguere peccantes},
to incite to sin and to accuse those who sin (Aleph 8; Apoc 12:10).

\paragraph{The laws of the angels.} On ``So shall I always keep thy law,
for ever and ever'' (Ps 118[119]:44), Hilary observes that there are laws of
the world to come as of this, and among them \latin{leges angelorum}: some
angels, standing with the Cherubim, praise God \latin{indefessa uoce},
with unwearied voice, while others stand before the face of God, unseen by
us, \latin{tamquam lege officii proprioris}, as by the law of a nearer
office; and the prophet knows \latin{pusillorum angelos deum cotidie ex
quadam lege conspicere}, that the angels of the little ones behold God
daily by a certain law (Vau 8).

\paragraph{The Church and its guard.} The heavenly Jerusalem, Hilary
teaches on Ps 124[125], is the Church, \latin{ecclesia angelorum
multitudinis frequentium}, crowded with the multitude of angels, the
Church of the first-born (Heb 12:22--23), and those who stand firm in it
lack neither the guard of the saints nor \latin{angelorum munitiones}, the
ramparts of the angels (124.4--5). The mountains round about Jerusalem are
the apostles, patriarchs, and prophets, or rather the angels, \latin{qui
ecclesiam quadam custodia circumsaepiunt}, who fence the Church about with
a kind of guard. Yet the psalm adds that the Lord is round about his
people, lest that guard be thought sufficient; Moses,
told that the angel would go before him, begged the Lord himself to go
with him (Exod 33:2, 15): \latin{bonum quidem praesidium angeli, sed melius
domini}---the guard of an angel is good, but the Lord's is better; for
the spiritual wickednesses and powers of this air fly about us on every
side (124.6).

\paragraph{The angels of churches, nations, and little ones.} On ``Let thy
ears be attentive to the voice of my supplication'' (Ps 129[130]:2),
Hilary explains that God has no bodily ears, and then turns to those
spirits who may be called his eyes and ears:
\begin{quote}
\latin{meminimus esse plures spiritales uirtutes, quibus angelorum est
nomen, uel ecclesiis praesidentes, sunt enim secundum Iohannem Asianis
ecclesiis angeli, sunt et Moyse testante secundum numerum angelorum fines
gentium Adae filiis constituti. sunt et domino docente pusillorum angeli
cotidie deum uidentes. sunt secundum Raphael ad Tobiam loquentem angeli
adsistentes ante claritatem dei et orationes deprecantium ad dominum
deferentes.} (\emph{Tract.\ in Ps.}\ 129.7)
\end{quote}
\noindent Angels preside over the churches of Asia (Apoc 1:20); the
bounds of the nations were fixed by the number of the angels (Deut 32:8,
in the Septuagint form Hilary reads); the angels of the little ones see
God daily; and Raphael tells of angels who stand before God's glory and
carry to him the prayers of those who pray (Tob 12:12, 15). Hilary draws the conclusion with care:
\latin{intercessione itaque horum non natura dei eget, sed infirmitas
nostra}. It is not God's nature that needs their intercession, but our
weakness, since they are ``ministering spirits, sent to minister for them
who shall receive the inheritance of salvation'' (Heb 1:14); God is
ignorant of nothing we do, but our infirmity needs the ministry of
spiritual intercession to ask and to obtain. Aquinas teaches the same when
he says that we pray to the saints, ``whether angels or men, not that God
may through them know our petitions, but that our prayers may be
effective through their prayers and merits,'' citing the angel with the
incense of the saints' prayers (\ST{II-II q.~83 a.~4}; Apoc 8:4).

\paragraph{Guardians of a weaker nature.} On Ps 134[135]:7, ``He bringeth
forth winds out of his stores,'' Hilary reads the winds as the angels of
whom it is written, ``Who makest thy angels spirits'' (Ps 103[104]:4), and
among them the angels of the
little ones who daily see God. They are sent for salvation, and our
weakness needed them:
\begin{quote}
\latin{hi igitur spiritus ad salutem humani generis emissi sunt; neque enim
infirmitas nostra nisi datis ad custodiam angelis tot tantisque
spiritalium caelestium nequitiis obsisteret, opus ad id fuit naturae
potioris auxilio.} (\emph{Tract.\ in Ps.}\ 134.17)
\end{quote}
\noindent These spirits were sent out for the salvation of the human race;
our weakness, unless angels had been given for its guard, could not have
withstood so many and so great spiritual wickednesses in high places;
there was need of the help of a stronger nature. So the Lord strengthened
the trembling Moses: ``Behold I will send my angel, who shall go before
thee'' (Exod 23:20, 23). Earlier in the same tractate Hilary names the
ranks to whom the Apostle's word about ``gods many'' (1~Cor 8:5) belongs
far more than to men: \latin{angelis, archangelis, thronis et
dominationibus, potestatibus, principatibus} (134.9); and he passes over
the wonder of \latin{thronorum, dominationum, principatuum et potestatum,
Cherubim et Seraphim, angelorum et archangelorum} to the works of God in
sky and sea (134.12).

\paragraph{Psalmody in the sight of the angels.} ``I will sing praise to
thee in the sight of the angels'' (Ps 137[138]:1) moves Hilary to a
full statement of the angels' presence to the faithful. The prophet
thinks it too little to sing before men, who see only bodily acts; he
hastens to do what spiritual natures behold, \latin{scit enim se sub
specula angelorum uitam omnem moresque agere, et ubique haec diuinorum
ministeriorum auxilia fidelibus cunctis adsistere}. He knows that he
passes his whole life and conduct under the watch of the angels, and that
everywhere these helps of the divine ministries stand by all the faithful,
as it is written, ``The angel of the Lord shall encamp round about them
that fear him'' (Ps 33:8[34:7]). The law was ordained by angels (Gal
3:19), and Eliseus showed that human weakness is defended by angels: when
the Syrian host surrounded the city, the prophet prayed that his servant's
eyes might be opened, ``and behold, the mountain was full of horses, and
chariots of fire round about Eliseus'' (4~Kgs 6:17). Before such angels
the psalmist sings, to please spiritual beings by the contemplation of
spiritual works (137.5).

\sectionguard
\subsection{Ambrose}\label{sec:fl-ambrose}

Ambrose of Milan (d.~397) writes of the angels chiefly as a theologian of
the Trinity. His treatises \emph{On the Holy Spirit} and \emph{On the
Faith} measure the angels against the Spirit and the Son, and in doing so
state what the angels are: creatures, limited, holy by grace, advancing
in knowledge, ministers.

\paragraph{Creatures circumscribed by their substance.} The Spirit is not a
creature, Ambrose argues, because He is infinite: ``Since then, every
creature is confined within certain limits of its own nature, and
inasmuch as those invisible operations, which cannot be circumscribed by
place and bounds, yet are closed in by the property of their own
substance; how can any one dare to call the Holy Spirit a creature, Who
has not a limited and circumscribed power?'' (\emph{De Spiritu Sancto}
I.7). The Spirit fills all things; ``Of what Angel does the Scripture say
this? of what Dominion? of what Power?~\dots\ For Angels were sent to few,
but the Holy Spirit was poured upon whole peoples.'' Aquinas cites this
chapter for the principle that spiritual creatures, though not in place
as bodies are, ``are none the less circumscribed by their substance''
(\ST{I q.~50 a.~1 obj.~3 and ad~3}; \S\ref{sec:q50}).

\paragraph{Holy by the grace of the Spirit.} The same chapter teaches that
the angels owe their holiness to the Spirit who sanctifies them:
\begin{quote}
But in like manner as the Spirit sanctifying the apostles is not a
partaker of human nature; so, too, He sanctifying Angels, Dominions, and
Powers, has no partnership with creatures. But if any think that the
holiness of the Angels is not spiritual, but some other kind of grace
belonging to the property of their nature, they will forsooth judge
Angels to be inferior to men.~\dots\ But since Angels come down to men to
assist them, it must be understood that the nature of Angels is higher as
it receives more of the grace of the Spirit, and that the favour awarded
to us and to them comes from the same author. (\emph{De Spir.\ S.}~I.7.83)
\end{quote}
\noindent ``But how great is that grace which makes even the lower nature
of the lot of men equal to the gifts received by Angels, as the Lord
Himself promised, saying: `Ye shall be as the Angels in heaven.' Nor is it
difficult, for He Who made those Angels in the Spirit will by the same
grace make men also equal to the Angels'' (I.7.84). Earlier Ambrose had
said that ``Whether you speak of Angels, or Dominions, or Powers, every
creature waits for the grace of the Holy Spirit,'' and that ``the nature
of Angels evidently can be changed'' (I.5.62--63). The angels'
holiness is grace, not nature; the higher nature receives the fuller
grace. Aquinas teaches that an angel needs grace to turn to God as the
object of beatitude, and that grace and glory were given to the angels
according to the degree of their natural gifts (\ST{I q.~62 aa.~2, 6};
\S\ref{sec:q62}).

The Spirit's river, Ambrose writes, ``seems to flow more abundantly among
those celestial Thrones, Dominions and Powers, Angels and Archangels,
rushing in the full course of the seven virtues of the Spirit,'' and makes
glad ``that heavenly nature of the creatures with the larger fertility of
His sanctification'' (I.16.178).

\paragraph{The Seraph a minister.} To those who cite the Seraph who purged
Isaiah's lips as proof that a creature forgives sins, Ambrose answers that
the coal from the altar is the grace of the Spirit, and that ``even if the
Seraph had taken away sin, it would have been as one of the ministers of
God appointed to this mystery. For thus said Isaiah: `For one of the
Seraphim was sent to me'\,'' (I.10.113--115; Isa 6:6--7). The contrast is
drawn in balanced clauses:
\begin{quote}
The Spirit, also, is indeed said to be sent, but the Seraph to one, the
Spirit to all. The Seraph is sent to minister, the Spirit works a
mystery. The Seraph performs what is commanded, the Spirit divides as He
wills. The Seraph passes from place to place, for he does not fill all
things, but is himself filled by the Spirit. (\emph{De Spir.\ S.}\
I.11.116)
\end{quote}

\paragraph{The angels before the mysteries.} In \emph{On the Faith} Ambrose
confesses that the generation of the Son is beyond the knowledge ``not
mine alone, but the angels' also. It is above Powers, above Angels, above
Cherubim, Seraphim, and all that has feeling and thought''; and he bids
the questioner, ``Do thou, then (like the angels), cover thy face with thy
hands, for it is not given thee to look into surpassing mysteries!''
(\emph{De fide} I.10.64--65; Isa 6:2). The angel's immortality is a gift:
``an angel is not absolutely immortal, his immortality depending on the
will of the Creator'' (III.3.19). At the Ascension the angels themselves
wondered:
\begin{quote}
For a Conqueror came, adorned with wondrous spoils, the Lord was in His
holy Temple, before Him went angels and archangels, marvelling at the
prey wrested from death, and though they knew that nothing can be added
to God from the flesh~\dots\ nevertheless, beholding the trophy of the
Cross~\dots\ they sought some broader and more lofty passage for Him on
His return. (\emph{De fide} IV.1.6)
\end{quote}
\noindent The angels bade their princes lift up the gates (Ps
23[24]:7--10); others asked ``Who is the King of glory?'' ``Howbeit,
seeing that the angels (as well as ourselves) acquire their knowledge
step by step, and are capable of advancement, they certainly must display
differences of power and understanding, for God alone is above and beyond
the limits imposed by gradual advance'' (IV.1.9--10). Aquinas teaches that
the angels learn the mysteries of grace by revelation, the higher more
deeply than the lower, and that even the higher learned particulars of
the Incarnation after the beginning (\ST{I q.~57 a.~5} and ad~1;
\S\ref{sec:q57}).

\paragraph{The angels to be entreated.} In his treatise \emph{On Widows}
Ambrose counsels the widow burdened with sin to seek intercessors, as
Peter and Andrew prayed for Peter's mother-in-law:
\begin{quote}
The flesh is weak, the soul is sick and hindered by the chains of sins,
and cannot direct its feeble steps to the throne of that physician. The
angels must be entreated for us, who have been to us as guards; the
martyrs must be entreated, whose patronage we seem to claim for ourselves
by the pledge as it were of their bodily remains. (\emph{De viduis}
9.55)
\end{quote}
\noindent The guardians given to us are also our intercessors, and the
faithful are to ask their prayers. The Church invokes the angels' help in
her liturgy (CCC~335), and Aquinas justifies prayer to angels and saints
as prayer that they intercede (\ST{II-II q.~83 a.~4}).

\sectionguard
\subsection{Jerome}\label{sec:fl-jerome}

Jerome (d.~420), translator and exegete, speaks of the angels in his
commentaries and letters as a reader of the Hebrew and Greek text. Three
of his judgements stand in the \emph{Summa} under his name: the guardian
angel of every soul from birth, the antiquity of the angels before our
world, and the princes of the nations in Daniel; his reading of the name
of the Seraphim agrees with the etymology Aquinas gives.\footnote{Jerome's
Latin follows Migne: Letter 18 in PL~22, the commentary on Daniel in
PL~25, and the commentaries on Matthew, Ephesians, and Titus in PL~26;
Letter 18 is cited by the sections of Vallarsi's text printed there.}

\paragraph{Each soul's angel from birth.} Commenting on ``See that you
despise not one of these little ones: for I say to you, that their angels
in heaven always see the face of my Father who is in heaven'' (Matt
18:10), Jerome reads the verse as a tempering of the severe precept that
precedes it. The Lord has commanded that hand, foot, and eye be cut off
when they scandalize, that is, every tie of kinship that leads to sin;
now he adds mercy to severity: as far as lies in you, despise no one, but
seek their healing by your own salvation; if they persist in sin, it is
better that you be saved alone than perish with many. Then the reason:
\begin{quote}
\latin{Quia angeli eorum in caelis vident semper faciem Patris. Magna
dignitas animarum, ut unaquaeque habeat ab ortu nativitatis in custodiam
sui angelum delegatum.} (\emph{In Matth.}\ III, on 18:10)
\end{quote}
\noindent Great is the dignity of souls, that each should have, from the
beginning of its birth, an angel delegated to its keeping. Jerome joins to
the verse the angels of the seven churches of the Apocalypse (Apoc
1:20--3:14) and the Apostle's command that women's heads be veiled in the
churches ``because of the angels'' (1~Cor 11:10). This sentence is the
\emph{sed contra} of Aquinas's article on the guardian of each man, and
decides his article on the moment the guardianship begins: the benefits
given to man as a Christian begin with baptism, those given to him as a
rational creature at birth; ``and with reason'' Jerome places the angel's
charge at birth (\ST{I q.~113 aa.~2, 4--5}; \S\ref{sec:q113}). The Church
teaches that ``from infancy to death human life is surrounded by their
watchful care and intercession'' (CCC~336).

\paragraph{The Seraphim (Letter 18).} Writing to Pope Damasus from
Constantinople in 381, Jerome expounds Isaiah's vision (Isa 6:1--7) and
asks what the Seraphim are, how they stand round God and fly, and why they
cry one to another. First the name: \latin{Seraphim sicut in
interpretatione Nominum Hebraeorum invenimus, ardor, aut incendium, aut
principium oris eorum, interpretantur}---the name means burning, fire, or
the beginning of their mouth (\emph{Ep.}\ 18.6). He finds the saving fire, which set the disciples' hearts burning on the
road to Emmaus (Luke 24:32), in the sacred books; and so the two
Seraphim, in his own mystical reading, are the two Testaments that stand
round God,
\latin{cum per ea Dominus ipse discatur}, since through them the Lord
himself is learned. Aquinas gives the same etymology when he explains why
the first angel who sinned is called a Cherub and not a Seraph: ``Seraphim
means `those who are on fire,' or
`who set on fire'\,'' (\ST{I q.~63 a.~7 ad~1}; cf.\ \ST{I q.~108 a.~5
ad~5}).

The wings that veil the face and feet are for Jerome the veils over God's
beginning and end: \latin{Velabant faciem non suam, sed Dei. Quis enim
ejus scire potest principium, quid antequam istum conderet mundum, in
rerum fuerit aeternitate: quando Thronos, Dominationes, Potestates,
Angelos, totumque ministerium caeleste condiderit?} They veiled not their
own face but God's; for who can know his beginning, what he was in the
eternity of things before he founded this world, when he made the
Thrones, Dominations, Powers, Angels, and the whole heavenly ministry?
The feet hide what will follow the consummation of the world; only the
middle, from creation to the Incarnation, Scripture opens to us (18.7).
He reports as well the reading of a Greek most learned in the Scriptures,
that the Seraphim are \latin{virtutes quasdam in caelis}, powers in
heaven who stand before God's judgement-seat and praise him, \latin{et in
diversa ministeria mittantur, maximeque ad eos qui purgatione indigent},
and are sent on divers ministries, above all to those who need purgation
(18.9). He closes on the daily grace of the vision: \latin{Quotidie
ad nos mittitur Seraphim, quotidie ingemiscentium atque dicentium: O miser
ego, quoniam compunctus sum, ora purgantur}. Every day a Seraph is sent to
us; every day the mouths of those who groan and say, Woe is me, for I am
pierced, are purged, and freed from sin they prepare themselves for the
ministry of God. And since the Septuagint makes \emph{Seraphim} neuter and
Symmachus masculine, he adds, \latin{Nec putandum sexum esse in Virtutibus
Dei}: there is no sex in the Powers of God (18.17).

\paragraph{The ages before our world.} The same question of the angels'
antiquity returns in Jerome's commentary on Titus. Expounding the life
God ``promised before the times of the world'' (Titus 1:2), he finds the
ages of the angels beyond reckoning: \latin{Sex mille necdum nostri orbis implentur anni: et
quantas prius aeternitates, quanta tempora, quantas saeculorum origines
fuisse arbitrandum est, in quibus angeli, throni, dominationes, caeteraeque
virtutes servierint Deo}. Six thousand years of our world are not yet
fulfilled; how many eternities, how many times, how many beginnings of
ages must we think there were before, in which the angels, thrones,
dominations, and the other powers served God (\emph{In Tit.}, on 1:2)!
Aquinas sets this text as the first objection in his article on the time
of the angels' creation and answers that ``Jerome is speaking according to
the teaching of the Greek Fathers''; he holds it more probable that the
angels were created with the corporeal world, and declares the contrary
opinion not erroneous (\ST{I q.~61 a.~3} and ad~1; \S\ref{sec:q61}).

\paragraph{The heavenly ministries.} On Eph 1:21, ``Above all principality
and power and virtue and dominion,'' Jerome asks where the Apostle found
these names, and answers that Paul, reading the Law spiritually, saw in
the kings, princes, captains, and centurions of Numbers and Kings the
image of other princes and kingdoms: \latin{quod scilicet in caelestibus
sint principatus, sint potestates, sint dominationes atque virtutes, et
caetera ministeriorum vocabula}, names of ministries that neither we can
name nor, he thinks, Paul himself, set in a heavy body, could number. An
earthly empire has its prefects, judges, and provinces, its counts,
dukes, and tribunes; \latin{et putamus Deum, Dominum dominorum et regem
regnantium, simplici tantum ministerio contentum?}---and shall we think
that God, the Lord of lords and King of kings, is content with a single
ministry (\emph{In Eph.}\ I, on 1:21)?

\paragraph{The princes of the nations (Daniel 10).} The angel who came to
Daniel said, ``the prince of the kingdom of the Persians resisted me one
and twenty days: and behold Michael, one of the chief princes, came to
help me'' (Dan 10:13). Jerome answers the question why the angel did not
come at once, though the prayer was heard on the first day, and takes the
prince of Persia for the angel entrusted with that nation:
\begin{quote}
\latin{Videtur mihi hic esse angelus cui Persis credita est, juxta illud
quod in Deuteronomio legimus: Quando dividebat Altissimus gentes et
disseminabat filios Adam, statuit terminos gentium juxta numerum
angelorum Dei.~\dots\ Restitit autem princeps, id est, angelus Persarum,
faciens pro credita sibi provincia, ne captivorum omnis populus
dimitteretur.} (\emph{In Dan.}, on 10:13)
\end{quote}
\noindent The prince is the angel to whom Persia was entrusted (Deut
32:8, again in the Septuagint form), resisting on behalf of the province
committed to him, lest the whole people of the captives be let go. He withstood
twenty-one days, \latin{enumerans peccata populi Judaeorum, quod juste
tenerentur captivi, et dimitti non deberent}, recounting the sins of the
Jewish people, that they were justly held captive and ought not to be
released. Jerome also
links these princes to the Apostle's ``princes of this world'' who knew
not the wisdom of God (1~Cor 2:6--8). Then Michael came, \latin{angelus
Michael qui praeest populo Israel}, the angel set over the people of
Israel; \latin{Principes autem primos, archangelos intelligimus}---by the
``chief princes'' we understand the archangels. At the angel's departure
``the prince of the Greeks'' entered before God \latin{ut accusaret et
Persarum principem atque Medorum}, that the kingdom of the Macedonians
might succeed; \latin{Et revera mira sacramenta Dei}, and truly God's
mysteries are wonderful (on 10:20--21). Aquinas reports Jerome's
exposition and Gregory's, and explains the princes' ``resistance'' as the
consulting of the divine will on the merits of the peoples; the
guardianship of nations he gives to the Principalities or perhaps the
Archangels, calling Michael ``one of the princes'' (\ST{I q.~113 aa.~3,
8}; \S\ref{sec:q113}). Cassian reads the prince of Persia otherwise
(\S\ref{sec:fl-cassian}).

\sectionguard
\subsection{John Cassian and Abbot Serenus}\label{sec:fl-cassian}

John Cassian, formed in the monastery at Bethlehem and in the desert of
Scete, completed his \emph{Conferences} between 426 and 428 to set before
the monks of Gaul the teaching of the Egyptian fathers. The two
conferences of Abbot Serenus (VII and VIII) give the desert's account of
the angels and demons in the spiritual combat. Serenus, a monk of
such purity that an angel was sent in a vision to take from him the
fiery humour of the flesh (VII.2), answers the questions of Germanus in
two night-long conversations.

\paragraph{Incitement, not compulsion.} Germanus complains that so great a
swarm of enemies drives the mind where it would not go. Serenus answers:
``we mean that they oppose our progress in such a way that we can think of
them as only inciting to evil things and not forcing.~\dots\ as there is
in them ample power of inciting, so in us there is a supply of power of
rejection, and of liberty of acquiescing.'' God's help fights on our side
``with much greater power than their hosts fight against us,'' and so ``no
one can be deceived by the devil but one who has chosen to yield to him
the consent of his own will'' (\emph{Conl.}\ VII.8).

\paragraph{How the demons touch the soul.} Spirit can be joined to spirit
imperceptibly, for there is ``some sort of similarity and kinship of
substance'' between them; ``But it is impossible for spirits to be
implanted in spirits inwardly or united with them in such a way that one
can hold the other; for this is the true prerogative of Deity alone''
(VII.10). The possessed are not a counter-example. The unclean spirit does
not enter ``the actual substance of the soul,'' but ``seizes on those
members in which the vigour of the soul resides, and laying on them an
enormous and intolerable weight overwhelms it with foulest darkness,'' as
wine or fever may do; this is why the Lord gave Satan power over Job's
flesh but said, ``only preserve his soul'' (VII.12; Job 2:6). Cassian
draws the principle that God alone penetrates spirits, and here states an
opinion of his own: though angels, archangels, and souls are spiritual
natures, ``yet we ought certainly not to consider them incorporeal. For
they have in their own fashion a body in which they exist, though it is
much finer than our bodies are~\dots\ from which it is clearly gathered
that there is nothing incorporeal but God alone'' (VII.13). Gennadius
repeats the doctrine (\S\ref{sec:fl-leo-gennadius}); Aquinas determines
that the angels are altogether incorporeal (\ST{I q.~50 a.~1};
\S\ref{sec:q50}).

\paragraph{The sand in the dark.} Do the demons then see our thoughts?
``Nobody doubts that unclean spirits can influence the character of our
thoughts,'' but only from without; ``they cannot possibly come near to
those which have not yet come forth from the inmost recesses of the
soul.'' They read the outward man: the monk who looks anxiously at the
sun and asks the hour has admitted gluttony; the one who groans silently
or changes colour has taken anger into his heart; ``And it is no wonder
that this is discovered by those powers of the air, when we see that even
clever men can often discover the state of the inner man from his mien
and look and external bearing'' (VII.15). Serenus gives an image from the
night:
\begin{quote}
For just as some thieves are in the habit of examining the concealed
treasures of the men in those houses which they mean to rob, and in the
dark shades of night sprinkle with careful hands little grains of sand
and discover the hidden treasures which they cannot see by the tinkling
sound with which they answer to the fall of the sand~\dots\ so these too,
in order to explore the treasures of our heart, scatter over us the sand
of certain evil suggestions, and when they see some bodily affection
arise corresponding to their character, they recognize as if by a sort of
tinkling sound proceeding from the inmost recesses, what it is that is
stored up in the secret chamber of the inner man. (\emph{Conl.}\ VII.16)
\end{quote}
\noindent Aquinas teaches the same distinction: a secret thought is known
to God in itself, but to angels and demons in its effects, ``not merely by
outward act, but also by change of countenance'' (\ST{I q.~57 a.~4};
\S\ref{sec:q57}).

\paragraph{Ranks of evil and the order of combat.} ``Not all devils can
implant all the passions in men, but~\dots\ certain spirits brood over
each sin,'' some over uncleanness, others over blasphemy, anger,
gloominess, or vainglory, ``and each one implants in the hearts of men
that sin, in which he himself revels'' (VII.17). There is no lasting
concord among wicked spirits, yet they arrange their assaults; they
``observe times and changes among themselves'' and haunt particular
places, and advance one by one, so that when one is beaten ``he must make
way for another spirit to attack it still more vehemently'' (VII.19). The
contest is proportioned: ``with beginners and feeble folk only the weaker
spirits join battle,'' and the assaults grow as the athlete of Christ
grows, for no saint could bear their malice ``were it not that the
merciful judge of our contest, and president of the games, Christ
Himself, equalized the strength of the combatants'' (VII.20; 1~Cor 10:13).
The demons labour and suffer confusion in defeat, and are destroyed ``with
a double destruction: first, because while men are seeking after
holiness, they, though they possessed it, lost it, and became the cause
of man's ruin; secondly, because being spiritual existences, they have
been vanquished by carnal and earthly ones'' (VII.21).

\paragraph{The limits God sets.} ``But that they have not the power of
hurting any man is shown in a very clear way by the instance of the
blessed Job, where the enemy did not venture to try him beyond what was
allowed to him by the Divine permission''; the demons themselves begged
leave to enter the swine, and could not of their own will enter even
unclean beasts, much less men made in the image of God (VII.22; Matt
8:31). The elders, Serenus adds, found the demons' power diminished since the first
anchorites, ``either the malice of the devils has been beaten back by the
power of the cross penetrating even to the desert, and by its grace which
shines everywhere; or else our carelessness makes them relax'' (VII.23).
They cannot enter a body until they have first taken possession of the
mind (VII.24), and those possessed by sins are more wretched than those
possessed by devils (VII.25). Therefore the possessed are not to be
despised, since ``none of them can be tempted at all by them without God's
permission,'' and all such trials are sent ``as by a most kind father and
most compassionate physician'' (VII.28). Nor were they kept from
communion: the elders ``thought that it ought to be given to them daily,''
for when received ``it, so to speak, burns out and puts to flight the
spirit which has its seat in his members'' (VII.30). Aquinas cites this
conference for giving Communion to those vexed by unclean spirits
(\ST{III q.~80 a.~9 ad~2}); the ordering of every assault to the good of
the elect he refers to God (\ST{I q.~114 a.~1}; \S\ref{sec:q114}).

\paragraph{The kinds of demons.} There exist ``in unclean spirits as many
desires as there are in men'': some only mock travellers, others delight
in bloodshed, others inspire pride or blasphemy, and Serenus reports ``a
demon openly confessing that he had proclaimed a wicked and impious
doctrine by the mouths of Arius and Eunomius.'' Scripture's names for
them, from asps and basilisks to ``the prince of this world,'' describe
``their fierceness and madness under the sign of those wild beasts''
(VII.32).

\paragraph{The fall of the angels.} In the second conference Germanus asks
whether the principalities and powers of Eph 6:12 were created to fight
against men. Serenus first sets down the rule: ``God forbid that we should
admit that God has created anything which is substantially evil, as
Scripture says `everything that God had made was very good'\,'' (VIII.6;
Gen 1:31). He then gives the tradition of the fathers on the creation of
the angels:
\begin{quote}
None of the faithful question the fact that before the formation of this
visible creation God made spiritual and celestial powers, in order that
owing to the very fact that they knew that they had been formed out of
nothing by the goodness of the Creator for such glory and bliss, they
might render to Him continual thanks and ceaselessly continue to praise
Him. (\emph{Conl.}\ VIII.7)
\end{quote}
\noindent Those ``present at the creation of the stars,'' who praised God
``with a loud voice'' (Job 38:7, Septuagint), were made before the heaven and the
earth, as Jerome had also taught, and as Aquinas grants may be held
without error (\ST{I q.~61 a.~3}); the Apostle enumerates them in order, ``whether they be angels or
archangels, whether they be thrones or dominions, whether they be
principalities or powers'' (Col 1:16 as Cassian cites it). Out of that
number ``some of the leaders fell,'' as Ezekiel's lament over the prince
of Tyre and Isaiah's over Lucifer show (Ezek 28:12--19; Isa 14:12--14);
the dragon drew with him ``the third part of the stars'' (Apoc 12:4); the
angels ``who kept not their first estate'' are reserved ``in everlasting
chains under darkness'' (Jude 6). Hence the ranks among the demons, for one of two
reasons: ``either that they still retain those differences of rank~\dots\
from the station of their former rank in which they were severally
created, or else that, though themselves cast down from heavenly places,
yet, as a reward for that wickedness of theirs in which they have
graduated in evil, they claim in perversity these grades and titles of
rank among themselves, by way of copying those virtues which have stood
firm there'' (VIII.8). Aquinas holds that the demons keep the orders of
their nature, which they have not lost (\ST{I q.~109 a.~1};
\S\ref{sec:q109}).

\paragraph{Pride, then envy.} Germanus had believed that the devil fell
through envy of Adam. Serenus shows from Genesis that the serpent was
already ``more subtle than all the beasts of the earth'' before the
temptation, a name no Scripture would give to Gabriel or Michael, and
concludes: ``that first fall of his, which was due to pride, and which
obtained for him the name of the serpent, was followed by a second owing
to envy,'' when he saw man ``called to that glory, from which he
remembered that he himself, while still one of the princes, had fallen''
(VIII.10). The order of the two sins is Aquinas's own (\ST{I q.~63 a.~2}).
The devil's lie and murder come from the same source: created ``a spirit
or an angel and good,'' with no Father but God, he ``became puffed up by
pride and had said in his heart: `I will ascend above the heights of the
clouds, I will be like the Most High,'\,'' and so became not only a liar
``but also the father of the actual lie'' (VIII.25; John 8:44).

\paragraph{The powers of the air and the princes of nations.} ``The
atmosphere which extends between heaven and earth is ever filled with a
thick crowd of spirits,'' and Providence has mercifully withdrawn them
from our sight, lest men be driven out of their wits by dread or corrupted
by their example (VIII.12).
Aquinas assigns the dark air to the demons as a place of punishment, that
they may tempt men for their exercise (\ST{I q.~64
a.~4}). Serenus reads Daniel's prince of Persia as ``a hostile power,
which favoured the nation of the Persians an enemy of God's people,'' and
the prince of the Greeks as another such patron; from which ``antagonistic
powers raise against each other those quarrels of nations,'' each
``striving with restless jealousy on behalf of those whom he presides
over, against the patron of some other nation'' (VIII.13). They are called
principalities and powers ``because they rule and preside over different
nations, and at least hold sway over inferior spirits and demons,'' and
they will be destroyed when Christ delivers the kingdom to the Father (1~Cor
 15:24; VIII.14). Aquinas reads the same verse of the ending of the
offices of leading others to the end (\ST{I q.~108 a.~7 ad~1}).

\paragraph{The names of the holy powers.} The same titles belong ``to the
better sort,'' and they are ``names of office and of worth or dignity'':
angels ``from their office of bearing messages,'' archangels ``because
they preside over angels, `dominions' because they hold dominion over
certain persons, and `principalities' because they have some to be
princes over, and `thrones' because they are so near to God and so privy
and close to Him that the Divine Majesty specially rests in them as in a
Divine throne'' (VIII.15). Aquinas gives the Thrones the immediate
knowledge of the types of the divine works (\ST{I q.~108 a.~5 ad~6}).

\paragraph{The tribunal of the demons.} That the demons obey worse powers
Serenus proves by Scripture and by a brother's vision. Sheltering in a
cave, the brother saw ``innumerable troops of demons'' gather round their
prince, who sat ``as on some lofty tribunal'' and examined each: the idle
he drove out with shame, the successful he dismissed ``with the highest
praise~\dots\ as most brave warriors''; one boasted that after fifteen
years he had that night overthrown a well-known monk, and the brother,
going to Pelusium, found the report true (VIII.16). Aquinas says that the
demons ``usurp a semblance of Divine power, by deputing certain ministers
to assail man, as the angels of God in their various offices minister to
man's salvation'' (\ST{I q.~114 a.~1}).

\paragraph{Two angels with every man.} ``Holy Scripture bears witness that
two angels, a good and a bad one, cling to each one of us.'' Of the good
the Saviour speaks in Matt 18:10; ``the angel of the Lord shall encamp
round about them that fear Him'' (Ps 33:8[34:7]); of Peter the disciples
said, ``It is his angel'' (Acts 12:15). ``But of both sorts the book of
the Shepherd teaches us very fully'' (VIII.17). The adversary who plotted
against Job ``always plotted against him but never could entice him to
sin,'' worsted ``by the Lord's protection which ever shielded him.''
Antony's victory over the demons sent against him by two philosophers
practised in magic confirms that
``demons cannot possibly find an entrance into the mind or body of
anyone, nor have they the power of overwhelming the soul of anyone,
unless they have first deprived it of all holy thoughts'' (VIII.19).

\paragraph{The sons of God.} Asked about the angels of Gen 6:2, Serenus
answers: ``We cannot possibly believe that spiritual existences can have
carnal intercourse with women.'' The sons of God are the line of Seth,
``by reason of its sanctity termed `angels of God,' or as some copies have
it `sons of God,'\,'' who took wives of the daughters of Cain and learned
their wickedness; the knowledge of nature handed down from Adam they
turned ``at the suggestion of devils to profane and harmful uses''; and
``in this way then that common notion, according to which men believe
that angels delivered to men enchantments and diverse arts, is in truth
fulfilled'' (VIII.21). Cassian thus keeps Tertullian's tradition of the
forbidden arts and gives it to men instructed by devils.

\sectionguard
\subsection{Leo the Great and Gennadius}\label{sec:fl-leo-gennadius}

\paragraph{Leo.} Leo the Great's letter to Turibius (447), with its
teaching \latin{diabolus bonus esset, si in eo quod factus est
permaneret}, belongs to the acts of the Church (DS~286; \S\ref{sec:faith}). His
sermons give the angels their place in the mysteries of Christ. At the
Ascension the apostles rejoiced, ``when in the sight of the holy
multitude, above the dignity of all heavenly creatures, the Nature of
mankind went up, to pass above the angels' ranks and to rise beyond the
archangels' heights, and to have Its uplifting limited by no elevation
until, received to sit with the Eternal Father, It should be associated
on the throne with His glory'' (\emph{Serm.}\ 73.4). The two angels at the
Ascension taught the Church ``not to doubt that all things are subjected
to Him on Whom the ministry of angels had waited from the first beginning
of His Birth''; as an angel announced the conception and heavenly voices
the birth, as messengers first attested the Resurrection, ``so the
service of angels was employed to foretell His coming in very Flesh to
judge the world, that we might understand what great powers will come
with Him as Judge, when such great ones ministered to Him even in being
judged'' (\emph{Serm.}\ 74.4). The Catechism summarizes the same life of
the Word incarnate as ``surrounded by the adoration and service of
angels'' (CCC~333). 

\paragraph{Gennadius, \emph{On the Dogmas of the Church}.} Gennadius, a
priest of Marseilles writing at the close of the fifth century, gathered
the Church's teaching into short articles, a book that was at times cited
under the name of Augustine; Aquinas cites it as \emph{De ecclesiasticis
dogmatibus} without an author.\footnote{Chapter numbers follow
Elmenhorst's text as printed in PL~58. The \emph{Summa} cites the same
chapters as xxix (ch.~59), xlix and lxxxii (ch.~82), and lxxxiii
(ch.~83).} Its articles on the angels are brief and exact.

The angels were made with the first creation: \latin{In principio creavit
Deus coelum, et terram, et aquam ex nihilo. Et cum adhuc tenebrae ipsam
aquam occultarent, et aqua terram absconderet, facti sunt angeli et omnes
coelestes virtutes, ut non esset otiosa Dei bonitas}. While darkness
still hid the waters, the angels and all the heavenly powers were made,
that God's goodness might not be idle but show itself in them through
long spaces before the visible world was formed and adorned (ch.~10).
With Cassian, Gennadius
confesses nothing incorporeal but God, and gives his reason:
\latin{Creatura omnis corporea est, angeli et omnes coelestes virtutes
corporeae, licet non carne subsistant. Ex eo autem corporeas esse credimus
intellectuales creaturas, quod localiter circumscribuntur}: the angels
are bodies, though not of flesh, because they are circumscribed in place,
like the soul and the demons (chs.~11--12); and they are immortal, having
no flesh by which to fall (ch.~13). Against Origen, no restoration awaits the demons: they will
not \latin{in angelicam qua creati sunt redeant dignitatem}, return to
the angelic dignity in which they were made, but the Lord's sentence stands, ``Depart
from me, you cursed, into everlasting fire, which was prepared for the
devil and his angels'' (ch.~9; Matt 25:41). The Council of Constantinople
of 543 anathematizes the same restoration (DS~411).

On the fall and the perseverance of the angels, Gennadius states the
doctrine that the schools received. Evil \latin{non est a Deo creata, sed
a diabolo inventa, quia et ipse bonus a Deo creatus est}: the devil,
created good and free, became the inventor of evil and envied in others
what he had lost in himself (ch.~57). The holy angels hold their good by fidelity, not by
nature:
\begin{quote}
\latin{Angeli vero, qui in illa qua creati sunt beatitudine perseverant,
non natura possident bonum, ut non mutarentur cum caeteris, sed arbitrio
servantes bonam voluntatem, et bonum conditionis, et fidem suo Domino
rependendo.} (\emph{De eccl.\ dogm.}\ 59)
\end{quote}
\noindent They keep a good will by their own choice, repaying their Lord
the good of their condition and their faith, and so are rightly called
holy angels. Chapter 60 repeats Leo's sentence
on the devil, \latin{bonus esset, si in eo quod factus est permaneret}.
The angels fixed in the love of God received, when the proud fell, the
reward that no rust of stealing fault can touch them, and they remain in
the contemplation of their Maker without end (ch.~61); after Satan's fall
with his legions, those who stood were so confirmed that they cannot fall
(ch.~62). Aquinas sets chapter 59 as the first objection of his article on
whether the angels were created in beatitude, and teaches that the
beatified angels cannot sin (\ST{I q.~62 a.~1 obj.~1}; a.~8;
\S\ref{sec:q62}).

Three articles on the demons and the mind close the book's angelology.
\latin{Internas animae cogitationes diabolum non videre certi sumus, sed
motibus eas corporis ab illo et affectionum indiciis colligi experimento
didicimus}: the devil does not see the soul's inner thoughts but gathers
them from the movements of the body and the signs of the affections; the
secrets of the heart are known to Him alone of whom it is written, ``thou only
knowest the heart of all the children of men'' (ch.~81; 3~Kgs 8:39).
\latin{Non omnes malae cogitationes nostrae semper diaboli instinctu
excitantur, sed aliquoties ex nostri arbitrii motu emergunt}; good
thoughts are always from God (ch.~82). And
the demons do not enter the soul in substance, but are joined to it by
pressure from without: \latin{Illabi autem menti illi soli possibile est
qui creavit}, to enter the mind is possible to Him alone who created it
(ch.~83). These are Cassian's conclusions in the form of articles.
Aquinas takes chapter 82 as the \emph{sed contra} of his article on
whether all sins come from the devil's temptation, and chapter 83 when he
denies that the devil dwells in Antichrist by indwelling, since ``the
Trinity alone dwells in the mind'' (\ST{I q.~114 a.~3}; \ST{III q.~8
a.~8 obj.~2 and ad~1}).

\sectionguard
\subsection{The Latin Fathers before Dionysius and Their
Reception}\label{sec:fl-reception}

Aquinas cites Ambrose, Jerome, Cassian, and the book of Gennadius by
name, and teaches much of the rest in his own words; where he determines
otherwise, his article stands against that Father's teaching. The councils
and the Catechism take up the goodness of the devil's created nature, the
angels' care from birth, and the prayer the Church makes with them and to
them.

\medskip
\begin{fourstable}{Witness}{Locus}{Teaching}{Received at}
Tertullian & \emph{Apol.}\ 22.3--8 & Fallen angels and their brood ruin
 men; subtlety and tenuity; ``every spirit is possessed of wings'';
 prophecy stolen, cures feigned & \ST{I q.~114 a.~4} (the demons' wonders
 are not miracles in the strict sense)\\
Tertullian & \emph{Apol.}\ 23.4, 15--16 & Demons confess under exorcism
 and depart at the name of Christ the Judge & ---\\
Tertullian & \emph{Adv.\ Marc.}\ II.10 & The devil made good, the wisest
 angel, corrupt by his own choice (Ezek 28) & DS~800; \ST{I q.~63 a.~7}
 and ad~1\\
Tertullian & \emph{De pat.}\ 5 & Impatience and envy toward man, God's
 image, the devil's first fault & \ST{I q.~63 a.~2} (pride first, envy
 after)\\
Tertullian & \emph{De carne Christi} 3, 6 & Angels assume true flesh
 from nothing; their nature spiritual yet bodily of its kind & \ST{I
 q.~51 aa.~1--2} (no natural body; assumed bodies of air)\\
Tertullian & \emph{De orat.}\ 3.3; 16; 29 & The angels' unceasing
 Sanctus; Christians ``candidates for angelhood''; the angel of prayer &
 CCC~335\\
Tertullian; Cyprian; Lactantius & \emph{De cultu fem.}\ I.2;
 \emph{De hab.\ virg.}\ 14; \emph{Div.\ inst.}\ II.14 & Angels fell
 through the daughters of men (Gen 6:2) and taught forbidden arts &
 Answered: Cassian, \emph{Conl.}\ VIII.21; \ST{I q.~51 a.~3 ad~6}\\
Minucius Felix & \emph{Octavius} 26--27 & Demons degraded from heavenly
 vigour; lurk in idols; confess and flee when adjured & ---\\
Cyprian & \emph{De zelo et livore} 4 & The devil fell by envy of man made
 in God's image & \ST{I q.~63 a.~2}\\
Lactantius & \emph{Div.\ inst.}\ II.8 & The devil a spirit made after the
 Son, fallen by envy of him & ---\\
Lactantius & \emph{Div.\ inst.}\ II.14.1; II.16 & Angels sent to
 guard mankind; they serve as officers of a province, wanting no honour &
 \ST{I q.~113 a.~1}; \ST{I q.~110 a.~1}\\
Hilary & \emph{Tract.\ in Ps.}\ 118 Aleph 7--8; 137.5 & The world full of
 angels; life lived under their watch & CCC~336\\
Hilary & \emph{Tract.\ in Ps.}\ 129.7; 134.17 & Angels of churches,
 nations, little ones; they carry prayers; guards of a weaker nature &
 \ST{I q.~113 aa.~1--3}; \ST{II-II q.~83 a.~4}\\
Ambrose & \emph{De Spir.\ S.}\ I.7 & Every creature circumscribed by its
 substance & \ST{I q.~50 a.~1 obj.~3 and ad~3}\\
Ambrose & \emph{De Spir.\ S.}\ I.7.83--84 & Angels holy by the Spirit's
 grace; the higher nature receives more grace & \ST{I q.~62 aa.~2, 6}\\
Ambrose & \emph{De fide} IV.1.9--10 & Angels wonder at the Ascension and
 advance in knowledge step by step & \ST{I q.~57 a.~5}\\
Ambrose & \emph{De viduis} 9.55 & The angels given as guards are to be
 entreated & \ST{II-II q.~83 a.~4}; CCC~335\\
Jerome & \emph{In Matth.}\ III (18:10) & Each soul has an angel from its
 birth & \ST{I q.~113 aa.~2, 4--5} (\emph{sed contra})\\
Jerome & \emph{Ep.}\ 18.6--7, 17 & Seraphim ``burning''; veiling God's
 beginning and end; a Seraph sent daily to purge & \ST{I q.~63 a.~7
 ad~1}; \ST{I q.~108 a.~5 ad~5}\\
Jerome & \emph{In Tit.}\ (1:2) & Angels served God for ages before our
 world & \ST{I q.~61 a.~3 obj.~1 and ad~1} (contrary more probable)\\
Jerome & \emph{In Dan.}\ (10:13, 20--21) & The prince of Persia the angel
 entrusted with Persia; Michael over Israel; princes are archangels &
 \ST{I q.~113 aa.~3, 8}\\
Cassian & \emph{Conl.}\ VII.10--13 & Demons press on the body, never
 enter the soul; only God incorporeal & \ST{I q.~50 a.~1} (angels wholly
 incorporeal)\\
Cassian & \emph{Conl.}\ VII.15--16 & Demons read thoughts by bodily
 signs, as thieves by sand & \ST{I q.~57 a.~4}\\
Cassian & \emph{Conl.}\ VII.17--23 & Spirits specialize; combat
 proportioned; no power without God's leave & \ST{I q.~114 a.~1}\\
Cassian & \emph{Conl.}\ VII.30 & Communion not refused to the possessed &
 \ST{III q.~80 a.~9 ad~2}\\
Cassian & \emph{Conl.}\ VIII.7--8 & Angels made before the visible world;
 leaders fell; demons keep or usurp ranks & \ST{I q.~61 a.~3}; \ST{I
 q.~109 a.~1}\\
Cassian & \emph{Conl.}\ VIII.10 & The devil's first fall by pride, the
 second by envy & \ST{I q.~63 a.~2}\\
Cassian & \emph{Conl.}\ VIII.12--14 & The air full of spirits; hostile
 patrons of nations; their principalities ended at the last & \ST{I q.~64
 a.~4}; \ST{I q.~108 a.~7 ad~1}\\
Cassian & \emph{Conl.}\ VIII.17 & A good and a bad angel with every man &
 \ST{I q.~113 a.~2}; \ST{I q.~114 a.~1}\\
Leo & \emph{Serm.}\ 73.4; 74.4 & Human nature raised above the angels;
 angels minister from the Nativity to the Judgment & CCC~333\\
Gennadius & \emph{De eccl.\ dogm.}\ 9 & No restoration of the demons &
 Constantinople (543), can.~9 (DS~411)\\
Gennadius & \emph{De eccl.\ dogm.}\ 59, 61--62 & The holy angels keep
 the good by will; confirmed after Satan's fall & \ST{I q.~62 a.~1
 obj.~1}; \ST{I q.~62 a.~8}\\
Gennadius & \emph{De eccl.\ dogm.}\ 81--83 & The devil does not see inner
 thoughts; not all evil thoughts are his; God alone enters the mind &
 \ST{I q.~57 a.~4}; \ST{I q.~114 a.~3}; \ST{III q.~8 a.~8}\\
\end{fourstable}
